%-------------------------------------------------------------------------------
%	ABSTRACT PAGE
%-------------------------------------------------------------------------------

\addchaptertocentry{Хураангуй}

\begin{center}
{\scshape\Large \univname\par}
{\scshape\large \facname\par}\vspace{0.5cm}
{\huge\textbf{{Хураангуй}} \par}
\bigskip
{\Large{\ttitle} \par}
\bigskip

{\normalsize \shortname \par}\addressname

\end{center}

\textit{\textbf{Түлхүүр үгс: \keywordnames}}
\bigskip

Байгууллагуудад 5S болон lean аргачлал нэвтрүүлэхэд шалгах хуудас цаасаар хөтлөгдөж, илэрсэн дутагдлыг хэн, хэзээ засахыг хянах, сарын тайланг гараар нэгтгэх зэрэг ажил их цаг зарцуулдаг бөгөөд мэдээлэл цаас, хүснэгт, чатад тарамдан шийдвэр гаргахад ашиглагддаггүй. Энэхүү дипломын ажлын хүрээнд байгууллагын өдөр тутмын ажил, 5S аудит, сайжруулалтын санаа, тайлагналыг нэг дор нэгтгэсэн вэб болон гар утасны платформыг зохион бүтээв.

Судалгааны хэсэгт 5S, Gemba, Kaizen, шатлалт өдөр тутмын хурал зэрэг бүтээмжийн аргачлалууд болон ижил төстэй таван программ хангамжийг шинжилж, системийн функциональ ба функциональ бус шаардлага, use-case загвар, класс, дараалал, үйл ажиллагааны диаграмм, давхаргат архитектур, өгөгдлийн сангийн схем, гол алгоритмуудын зохиомжийг боловсруулав. Систем нь NestJS, PostgreSQL дээр суурилсан REST API, React вэб апп, сүлжээгүй орчинд ажилладаг Flutter гар утасны аппаас бүрдэнэ. Аудитын оноо босгоос доош бол засах ажлыг талбайн хариуцагчид автоматаар оноох, сарын тайланг бүртгэлээс автоматаар бүрдүүлж архивлах, сүлжээгүй үед хийсэн бүртгэлийг давхардуулалгүйгээр дараа нь илгээх зэрэг нь системийн шинэлэг тал юм.
